% !TEX root = ../main.tex

% 分析与讨论	
\begin{table}
	\renewcommand\arraystretch{1.7}
	\begin{tabularx}{\textwidth}{|X|X|X|X|}
		\hline
		Major：& Physics &Grade：& 2022\\
		\hline
		Name： & 杨舒云 & Student number：& 22344020 \\
		\hline
		Date：& 2024/9/20 & Score： &\\
		\hline
	\end{tabularx}
\end{table}
\section{锁相放大器工作原理、基本参数与基本操作  \\ Analysis \& Discussion}


%---------------------------------------------------------------------
% 数据处理
\subsection{Data Processing}

% 分析
\subsubsection{Analysis}
分析部分主要基于实验讲义中所提出的问题做分析。
\begin{ubox}{}
	【问题 1】 示波器是否显示出正弦波？其测量的信号参数值是否与锁相放大器的信号设置值一致？
	
	【问题 2】 请比较示波器的读数和锁相放大器的$R$值，以理解$R$值的确切含义。
\end{ubox}
由\cref{fig:D1-1-1}可知，示波器显示的是正弦波，测量值为 $V_{pp} = 240.0 \mathrm{mV}$，因此测量的有效值为
\[
V_{rms} = \frac{V_{pp}}{2 \sqrt{2}} = \frac{240.0 \mathrm{mV}}{2 \sqrt{2}} \approx 84.85 \mathrm{mV}
\]

同时，锁相放大器测量的$R = 81.42 \mathrm{mV}$，与示波器测量的有效值相近，说明锁相放大器测量的$R$值为信号的有效值。

锁相放大器的设置值$V_{set} = 80 \mathrm{mV}$，与测量值有一定差异，但在误差范围内可认为是一致的。

此外，测量的频率$f = 1.000 \mathrm{kHz}$，与设定值一致。

\begin{ubox}{}
	【问题 3】 测量数据反映出$R$、$X$、$Y$及$\theta$的之间是什么关系？是否与式(D1- 22)、式(D1- 25)和(D1- 26)符合？
	
	【问题 4】 你使用的 OE1022 是双锁相放大器吗？为什么？
\end{ubox}
根据\cref{tbl:D1-2-1}的数据，我们可以得到：
\[
R \cos\theta = 81.43 \mathrm{mV} \cdot \cos(177.33°) = -81.342 \mathrm{mV} \approx X 
\]
\[
R \sin\theta = 81.43 \mathrm{mV} \cdot \sin(177.33°) = 3.793 \mathrm{mV} \approx Y 
\]
\[
\sqrt{X^2 + Y^2} = \sqrt{(-81.34 \mathrm{mV})^2 + (3.79 \mathrm{mV})^2} = 81.428 \mathrm{mV} \approx R
\]

因此是符合的。

同时，由上述关系可知，该锁相放大器可同时测量完整的输入信号信息（即信号的有效值$R$、相位角$\theta$和频率），因此我们使用的是双锁相放大器。

\begin{ubox}{}
	【问题 5】 示波器上看到什么波形？$X$、$Y$是正弦波吗？其相位差为多少度？$X$、$Y$的平均值为零吗？平均值是否会随滤波器参数而改变？参照式(D1-19)和式(D1-23)讨论；$R$值波形是否呈现全波整流的正弦信号，其频率为多少？与$X$或$Y$的频率存在什么关系？请对比式(D1- 28)讨论此波形。
\end{ubox}
首先，从公式考虑：
\begin{align*}
	u_{px}(t) &= \frac{1}{2} A_I s(t) \{ [ \cos((\omega_m - \omega_r)t + \theta) - \cos((\omega_m + \omega_r)t + \theta) + n(t)\sin(\omega_r t) ]\} \\
	&\approx \frac{1}{2} A_I s(t)\{ \cos\theta - \cos(2\omega_m t + \theta) + n(t)\sin(\omega_m)t\} \\
	u_{py}(t) &= \frac{1}{2} A_I s(t) [ \sin((\omega_m - \omega_r)t + \theta) + \sin((\omega_m + \omega_r)t + \theta) ]
\end{align*}

信号分为$\cos\theta$直流项，$\cos(2\omega_m t + \theta)$正弦波项。正如我们所观察\cref{fig:D1-3-3}可知，$X$、$Y$的波形并不是纯正的正弦波，而是正弦波加上了一个直流偏置，且由\cref{fig:APL1_1_question5_phase}知两者之间相位相差接近90°。

由\cref{tbl:D1-3-1}知，平均值不为0，会受滤波器参数的影响（考虑$u_{ox}(t)=A_I[\frac12s(t)\cos\theta+n_x(t)]$和$u_{oy}(t)=\frac12A_Is(t)\sin\theta$）。可能由于滤波器带宽调整导致信号频谱的变化，增益设置影响了信号的幅度，特别是在信号较弱或噪声较大的情况下从而影响了滤波效果，故平均值受滤波器参数影响。

关于滤波器参数的问题，详见后续几个问题。

理论上计算$R=\frac{1 - \cos(2\omega_m t + 2\theta)}{2}$，因此输出的双倍频正弦波类似于全波整流信号的形态，因为该信号的频率是输入信号的 2 倍。与输入信号相比，它的频率加倍。这是因为在乘法器中，输入信号与参考信号相乘会产生二次谐波成分，即频率翻倍的分量。如\cref{fig:D1-3-1}所示。

此外，也可以看出R与X、Y波形的频率相同。

\begin{ubox}{}
	【问题 6】 当调大时间常数后，输出$R$是否越来越稳定？直流分量（信号平均值）是否有改变？
	
	【问题 7】 观察示波器波形，并给出在某一陡降值下，时间常数为多少个信号周期时，才能观察到稳定的直流输出？提示：注意调节示波器$Y$轴量程以提高测量振幅变化的精度。并结合式(D1-32)讨论。
	
	【问题 8】为何在 1 阶（6dB/oct）时还没怎么起滤波效果，而到 4 阶（24dB/oct）时，则可看见明显的效果？
\end{ubox}

滤波器的参数（如带宽）会影响滤波效果，带宽越窄，滤波器对高频分量的抑制效果越好，输出波形的平均值越接近纯粹的直流分量。根据\cref{fig:apl1_question7}，调大时间常数后，能够更有效地抑制高频噪声和不相关的频率成分，从而让R信号中的高频成分被滤除，信号平均值越稳定。由\cref{tbl:D1-3-1}可知，增大时间常数（即改变了滤波器参数）,平均值有一定变化。因为由于高频分量的存在，波形的平均值不为零。

同时，我们也可发现在10ms时就已经有稳定的直流输出了。关于据式(D1-32)的理论分析，详见\textbf{原理概述}部分。

使用  $10\mu s$ 的时间常数，选择时间常数 \( TC \) 等效于选择滤波器的电阻 \( R \) 或电容 \( C \)，决定了一阶滤波器的带宽；而选择陡降值则等于选择级联 RC 滤波器的阶数 \( n \)，这决定了带通边缘的陡度。

观察\cref{fig:D1-4-1-1}、\cref{fig:D1-4-1-2}、\cref{fig:D1-4-1-3}和\cref{fig:D1-4-1-4}，可以发现，陡降设置的越大，输出波形的幅值越小，滤波效果越明显。对于陡降（Slope），陡降表示信号在截止频率处的衰减速度，滤波器的阶数越高，陡降越陡，能更有效地去除高频噪声。所以，陡降越陡，滤波器在其截止频率附近的频率选择性就越好，即滤波效果越明显。在 1 阶（6 dB/oct）时，滤波器对高频信号的抑制效果相对较弱，因为其带宽较宽，许多高频成分仍然可以通过。而当提升至 4 阶（24 dB/oct）时，滤波器的陡降增加，使得高频信号在截止频率附近的衰减更加显著。因此，在这种情况下，更多的高频成分被有效削减，输出信号的稳定性和直流分量的增强更加明显。这解释了为何在低阶滤波时滤波效果不明显，而高阶滤波时则能显著改善信号质量。


\begin{ubox}{}
	【问题 12】如何获得幅值为 0.1mVrms 的正弦波输入信号？
	
	【问题 13】对低信噪比信号或非周期噪声测量，能否用示波器的“auto”键？若通过手动调节，示波器屏宽（时间轴范围）应为多少个信号周期为宜？
	
	【问题 14】在用示波器观察、测量信号有效值时，请注意屏幕时间轴范围应达到待观测信号完整周期的整数倍；为什么？
	
	【问题 15】用锁相放大器测量 $R$ 值时，读数应保留几位有效数字？调节锁相放大器的什么参数设置可以得到？调节锁相放大器的什么参数可以获得稳定的读数。
\end{ubox}
由于信号发生器最低只能设置 0.7mVrms 的正弦波输入信号，所以使用信号发生器产生一个 1 Vrms 的正弦波信号，再经过噪声发生器的信号端接口的 -80dB 的衰减，即可得到幅值为 0.1mVrms 的正弦波输入信号。
\[
- 80 \mathrm{dB} = 20 \log \frac{V_{out}}{V_{in}} = 20 \log \frac{0.1 \mathrm{mVrms}}{V_{in}}\Rightarrow V_{in} = 1 \mathrm{Vrms}
\]

示波器的“auto”模式会首先检测输入信号的特征，包括频率、幅度和波形类型，根据信号的周期性，示波器会自动选择合适的时间轴范围，以确保信号在屏幕上完整显示；而非周期信号本身缺乏稳定的重复模式，因此很难通过自动设置来捕捉到完整的信号特征。因此，使用“auto”模式会导致采样过少，无法准确计算一个稳定的噪声幅值。手动调节时，不妨大胆将屏宽拉大，知道出现稳定幅值时就停止。

信号有效值的计算基于信号的平方在时间上的平均$u_{rms} = \frac{\int_0^{t} (u(\tau))^2 \mathrm{d\tau}}{t} \approx \frac{\int_0^{T} (u(\tau))^2 \mathrm{d\tau}}{T}$仅在“积分时间$t$远大于周期$T$”，或“积分时间$t$是周期$T$的整数倍”时才成立。所以，屏幕时间轴范围应达到待观测信号完整周期的整数倍。

读数应在取得2-3位稳定示数后，保留2-3位有效数字。通过设置合适的Sensitivity，可以使得有效数字位数显示的尽可能多且保证一定的测量精度。通过设置合适的时间常数和陡降，可以获得稳定的读数。


% 讨论
\subsubsection{Discussion}
讨论部分主要讨论实验过程中我们遇到的一些有意思的情况。
\begin{enumerate}
	\item 首先是关于实验1.3中我们观察到了异常波形，如\cref{fig:APL1_1_question5_wrong(AC)}所示。
	
	\begin{figure}[htbp]
		\centering
		\includegraphics[width=0.45\textwidth]{APL1_1_question5_wrong(AC).jpg}
		\caption{交流耦合下的乘法器异常波形}
		\label{fig:APL1_1_question5_wrong(AC)}
	\end{figure}
	
	通过老师的讲解，我们意识到它的产生是因为输出中产生了直流偏置，即$R=|A\sin(\omega t + \phi) + bias|$。事实上，在\textbf{思考题1.3}中我们就已经考虑过相关问题：“在实际测量中，直流偏置可能会引入仪器的零点漂移或偏置。然而，锁相放大器通常具有直流耦合和交流耦合模式，以及直流偏置调零功能，可以消除这些影响。锁相放大器通过同步检测，仅提取与参考信号频率和相位匹配的信号分量。其他频率（包括直流）和相位的信号分量会被滤除或平均为零”。
	
	进一步排查，我们的耦合方式（Coupling）选择的是交流耦合（AC），将它改为直流耦合（DC）后该异常波形消失，出现正常波形，如\cref{fig:APL1_1_question5_correct(DC)}所示。
	
	\begin{figure}[htbp]
		\centering
		\includegraphics[width=0.45\textwidth]{APL1_1_question5_correct(DC).jpg}
		\caption{直流耦合下的乘法器正常波形}
		\label{fig:APL1_1_question5_correct(DC)}
	\end{figure}
	
	\item 其次，在完成1.8时，黄罗琳同学（22344001）的建议非常对我们非常有启发，即在完成这个实验时，我们首先要熟悉其中的数据的概念：
	\begin{itemize}
		\item 信噪比公式，这是计算的基础：$\mathrm{SNR_{i}(dB)}=20lg(V_{S,RMS}/V_{N,RMS})$
		\item 正弦波 $V_{in}$幅值：这个需要根据，设定的信噪比以及实际的噪声信号的大小而定。
		\item 噪声信号大小：由于存在仪器老化，尽管教学实验箱噪声发生器的标准发出的噪声为$100mV_{rms}$，但是实际上存在偏差，需要测出一个实际值记录。
		\item $\mathrm{SNR_i}$：此信噪比相比于设定的输入信号信噪比会存在些许的不同，这是由于实验设定值会存在一定的误差，故存在不同。
		
		\item 滤波器带宽：带宽的计算是在设置示波器的数字滤波是，会给出一个频率上限和下限，将两者相减便可以得到带宽。
		\item 滤波后的信号有效值：在示波器上，使用滤波后会有一条紫色的线，这就是滤波后的信号，测得此信号的有效值即可。
		\item $SNR_{o,os}$：示波器滤波器滤波后的信号的和噪声有一个新的信噪比。
		\item 量程灵敏度：通过设置量程灵敏度，可以调整数值的精度，所以选择合适的灵敏度，可以既保证了量程满足输出的数值要求，也保证了一定的精度。
		\item 时间常数：由设置的频率（带宽）计算而来。
		\item 陡降：同样为预先设置Slope。
		\item LPF带宽：需要通过计算得来，根据不同的陡降得到不同的系数（来自讲义的表），在乘以时间常数的倒数便得到带宽。
		\item 信号有效值R：可以通过锁相放大器面板直接输出。
		\item 噪声测量值N：可以通过锁相放大器面板直接输出。
		\item  $SNR_{o,ol}$：锁相放大器滤波后的信号的和噪声有一个新的信噪比
	\end{itemize}
	
\end{enumerate}

% 总结
%\subsubsection{Conclusion}
%\begin{itemize}
%	\item 
%\end{itemize}


%---------------------------------------------------------------------
% 实验后思考题
\subsection{Reflections after Experiment}

% 思考题1
\begin{question}
	描述用教学实验箱配制不同信噪比信号的原理；实验箱生成的噪声是否为白噪声？
\end{question}
在实验教学中，生成不同信噪比（SNR, Signal-to-Noise Ratio）信号的目的是研究噪声对信号测量和处理的影响。教学实验箱通过控制信号与噪声的相对幅度来实现不同信噪比的信号生成。实验箱通常具有一个信号发生器，可以生成不同类型的周期性信号，如正弦波、方波、三角波等。这些信号代表我们希望测量或分析的主要目标信号。实验箱中也有一个噪声发生模块，用于生成随机噪声（如白噪声）。这些噪声是为了模拟真实环境中信号所受到的随机干扰。将信号与噪声进行线性叠加，即把信号和噪声的电压信号相加，得到一个包含噪声的信号，其数学表达式为：

\[
V_{\text{out}}(t) = V_{\text{signal}}(t) + V_{\text{noise}}(t)
\]

为了配制不同信噪比的信号，可以通过调节信号发生器的输出幅度来控制信号在总输出信号中的权重，从而改变信噪比。同样地，也可以通过调节噪声发生器的输出幅度来控制噪声对信号的干扰程度。如果噪声幅度增加，信噪比会降低，信号受到的干扰相对增加；如果噪声幅度减小，则信噪比会提高。信噪比通常定义为信号功率与噪声功率的比值，通常用分贝（dB）表示，其公式为：

\[
\text{SNR (dB)} = 10 \log_{10}\left(\frac{P_{\text{signal}}}{P_{\text{noise}}}\right)
\]

通过调节信号和噪声的幅度，可以改变信号和噪声的功率比值，从而得到不同的信噪比。通过实验箱生成不同信噪比的信号，可以研究在高噪声环境下信号处理的方法和技巧，例如使用滤波器来降低噪声对信号的影响，用锁相放大器提取特定频率信号中的有效成分，或比较不同信噪比下信号测量的精度和可靠性。

白噪声是一种频率谱均匀分布的随机噪声，这意味着它在所有频率上都具有相同的功率谱密度。理想的白噪声信号在理论上具有无限的带宽，能量均匀分布在整个频率范围内。教学实验箱生成的噪声通常是“白噪声”的一种近似实现。由于硬件电路和信号发生器的频率响应存在带宽限制，实际生成的白噪声并不是理想的无穷带宽，而是有限带宽内近似均匀的频谱。实验箱的白噪声频谱一般在某个频率范围内近似平坦，通常这部分频率是设备的有效频段，例如从几十赫兹到几十千赫兹或更高。在这个有效频带内，噪声可以认为是白噪声，具有均匀的功率谱密度。由于电子设备的硬件限制，实验箱生成的白噪声在频率响应上会存在一定的衰减，尤其是在较高频率和较低频率时。因此，教学实验箱产生的噪声只能称作“有限带宽的白噪声”或“准白噪声”。

白噪声在实验中的作用主要是用于模拟真实环境中的随机噪声。白噪声可以很好地模拟真实信号测量中随机存在的干扰，帮助学生理解和解决在实际测量中如何应对噪声问题。另外，白噪声也可以用于滤波器性能测试，当白噪声作为输入信号通过滤波器时，可以观察滤波器对不同频段的响应，从而判断滤波器的性能。通过对不同信噪比信号的实验，可以了解噪声对信号测量的影响，以及如何使用滤波器、锁相放大器等工具来提取有效信号，抑制噪声干扰。这些实验对于掌握信号与噪声的基本概念，以及理解信号处理技术的实际应用非常有帮助。

% 思考题2
\begin{question}
	随着信噪比的改变，信号 $R$ 和相位差$\theta$的测量值会有什么影响？
\end{question}
信噪比（SNR）对信号幅度 \( R \) 和相位差 \( \theta \) 的测量有着直接的影响，尤其是在噪声逐渐增大、信噪比降低的情况下，信号的测量精度和稳定性都会受到不同程度的干扰。为了深入理解这一影响，可以结合公式进行分析。

对于锁相放大器，其输出信号包括两部分：同相分量 \( X \) 和正交分量 \( Y \)。通过这两个分量，可以得到信号的幅度 \( R \) 和相位差 \( \theta \)：

\[
X = R_{\text{signal}} \cos(\Delta \phi)
\]
\[
Y = R_{\text{signal}} \sin(\Delta \phi)
\]

这里，\( R_{\text{signal}} \) 表示信号的真实幅度，\( \Delta \phi \) 表示待测信号与参考信号之间的相位差。信号的幅度 \( R \) 可以通过 \( X \) 和 \( Y \) 计算得出：

\[
R = \sqrt{X^2 + Y^2}
\]

信号的相位差 \( \theta \) 可以通过以下公式计算得到：

\[
\theta = \text{atan2}(Y, X)
\]

噪声对信号幅度的影响主要体现在信号的同相分量和正交分量中。热噪声和散粒噪声通常表现为随机的电压波动，且这些波动会加在信号的同相和正交分量上。对于幅度 \( R \) 的计算：

\[
R = \sqrt{(X_{\text{signal}} + X_{\text{noise}})^2 + (Y_{\text{signal}} + Y_{\text{noise}})^2}
\]

当信噪比较高时，噪声分量 \( X_{\text{noise}} \) 和 \( Y_{\text{noise}} \) 较小，对信号的总幅度 \( R \) 的影响不明显，测量值与真实值接近。然而，当信噪比降低时，噪声分量可能会变得显著，导致测量的信号幅度波动增加，幅度的精度降低。在低频测量中，1/f 噪声（闪烁噪声）会对信号的低频部分造成显著影响。如果信号本身的频率较低，那么随着信噪比的降低，1/f 噪声的影响会显得更加突出。这会导致信号幅度在不同的测量时间点之间产生明显的波动，从而影响测量的稳定性。

相位差的测量更为敏感，因为噪声的相对影响会显著干扰相位计算，尤其是当信号幅度较低时。相位差的测量公式为：

\[
\theta = \text{atan2}(Y, X)
\]

噪声对 \( X \) 和 \( Y \) 的干扰会直接影响相位的计算。当噪声幅度增大时，测量得到的同相分量 \( X \) 和正交分量 \( Y \) 都可能与实际值出现较大偏差，这种偏差会导致相位差 \( \theta \) 的测量出现不稳定。当信噪比降低时，信号被噪声所掩盖，特别是在信号幅度与噪声幅度接近时，相位测量结果会受到很大影响，甚至可能随机地出现大幅波动。这种波动通常表现为：当 \( X \) 或 \( Y \) 非常小，噪声对它们的影响相对更大，这时小的噪声会显著改变相位角的方向，从而使相位测量变得非常不稳定。当噪声非常强时，可能会出现“相位跳跃”（Phase Jitter），即相位测量值在多个周期内突然大幅变化，影响信号的稳定性和测量的可靠性。如果信号中存在高频噪声，这些噪声在经过锁相放大器的低通滤波器时可能部分被衰减，但剩余的高频噪声仍可能影响相位测量的准确性，尤其是在噪声的频率接近参考信号频率时，噪声的叠加可能导致相位偏移。

随着信噪比的降低，信号幅度 \( R \) 的测量精度和稳定性会逐渐降低。噪声分量的增加会导致信号幅度的波动增大，难以获得稳定的测量值。此外，相位差 \( \theta \) 的测量受噪声影响更加明显，尤其是在噪声较强时，相位的随机波动和跳跃现象会显著增加。对于低信噪比的情况下，相位测量可能变得完全不可靠，因为噪声的影响会使得同相分量和正交分量的计算结果发生较大偏差。因此，为了准确测量信号的幅度和相位，保持较高的信噪比非常重要。在实验中，可以通过增大信号幅度、减小噪声源的影响、使用滤波器减少高频或低频噪声等方法来提高信噪比，从而提高信号幅度和相位测量的精度和稳定性。

% 思考题3
\begin{question}
	所有的噪声都可以用锁相放大器消除吗？锁相放大器的理论检测极限是多少？受什么限制？
\end{question}
锁相放大器的核心功能是通过参考信号与输入信号相乘，并通过低通滤波器提取目标频率的信号成分，这使得它在处理与参考信号频率不同的噪声时非常有效。然而，并不是所有的噪声都可以被完全消除。对于带外噪声，锁相放大器对其抑制效果非常好。带外噪声是指频率远离参考信号频率的噪声，例如随机的高频或低频噪声，由于它们在频域上与目标信号不重叠，经过锁相放大器的处理后会被有效抑制。然而，当噪声的频率与目标信号的参考频率相同（例如电源噪声 50 Hz 的干扰与信号频率相同），这种同频噪声难以通过锁相放大器消除。这种噪声与信号频率一致，且可能具有随机相位，因此它也会出现在锁相放大器的输出中，导致检测精度受到严重影响。此外，锁相放大器对相位不稳定的噪声抑制效果也有限。例如，相位抖动或相位噪声直接影响信号的相位测量。此外，1/f 噪声（闪烁噪声）在低频范围内表现显著，尤其在测量低频信号时，其影响尤为突出。锁相放大器只能对目标频率上较大幅度的噪声进行抑制，对于频谱在低频范围内逐渐升高的 1/f 噪声，其抑制效果有限。因此，1/f 噪声的影响在低频信号测量中仍然是一个主要的问题。

锁相放大器的理论检测极限可以理解为它能够从多么低的信噪比环境中提取出有效的信号。这个检测极限受到多种因素的限制，包括噪声、放大器自身的特性、以及实验条件。通常情况下，锁相放大器能够将信噪比非常低的信号提取出来，其理论检测极限由以下几个方面来决定。首先，系统噪声限制会对测量的最低可检测信号产生影响。仪器内部的电子元件都会产生噪声（如热噪声和散粒噪声），这些噪声会限制锁相放大器的最小可检测信号。仪器的前端放大器部分的输入噪声水平越低，锁相放大器的检测极限就越高。此外，锁相放大器依赖于参考信号来进行相敏检测，因此参考信号的稳定性直接影响检测极限。如果参考信号存在相位漂移或抖动，锁相放大器的同步检测效果会受到影响，从而增加测量的不确定性和噪声，降低检测的灵敏度。

锁相放大器的测量带宽（即低通滤波器的时间常数）也会影响其检测极限。带宽越窄，锁相放大器可以更有效地抑制噪声，但同时它也会限制响应速度。因此，选择合适的带宽是一个权衡灵敏度和响应速度的过程。理论上，如果将带宽无限减小，则可以消除绝大部分噪声，但是这样也会导致信号响应变得极其缓慢。积分时间或低通滤波器的时间常数越长，测量中的随机噪声影响就会越小，信号的检测极限就越低。长时间的积分可以有效提高信号的信噪比，但也会导致系统响应变慢。因此，锁相放大器的极限检测能力也与积分时间密切相关。

理论上，锁相放大器的检测极限是可以非常低的，甚至可以在极低信噪比的环境下仍能提取出目标信号。然而，实际中的检测极限受限于系统本身的噪声水平、参考信号的稳定性、测量带宽和积分时间等因素。此外，锁相放大器对不同类型的噪声抑制效果不尽相同，对于与信号频率相同的噪声和相位相干噪声，锁相放大器的检测能力会显著降低。因此，尽管锁相放大器在抑制带外噪声方面非常有效，但并不能消除所有类型的噪声，尤其是频率和相位与信号一致的噪声。提高锁相放大器的检测极限需要从优化系统设计、提高参考信号稳定性、适当调整带宽和积分时间等方面入手。

% 思考题4
\begin{question}
	请比较锁相放大器与示波器（经过带通滤波后）测量不同信噪比信号的结果，分析两种测量仪器（方法）的优势与劣势。选择锁相放大器窄带宽的 LPF 可以降低式(D1- 34)中的$n_{\omega_r}$，从而提高测量值的信噪比吗？信噪比的提高与 LPF 带宽有何关系？参考式(D1-10)讨论。
\end{question}
\cref{tbl:D1-8-1}中展示了示波器和锁相放大器在三种不同输入信噪比（0 dB, -30 dB, -80 dB）情况下测量信号的结果。示波器在测量信号时，测量的有效信号幅值、噪声幅值以及经过带通滤波器处理后的输出信号与输出信噪比（SNRo,os）均列在表中。从表格中可以看出，在 0 dB SNR 输入信号下，输入信号幅值为 119.2 mVrms，噪声幅值为 120.7 mVrms，信号和噪声基本相当。经过带通滤波器后，输出信号幅值为 135 mVrms，输出信噪比 \( SNR_{o,os} \) 提高到 0.97 dB。在 -30 dB SNR 输入信号下，输入信号幅值为 3.8 mVrms，噪声幅值依然是 120.7 mVrms。在这种情况下，噪声远大于信号，经过带通滤波后输出信号幅值为 4.8 mVrms，输出信噪比为 -28.01 dB，与输入相比略有改善。而在 -80 dB SNR 输入信号下，示波器无法准确测量信号（标记为“-”），说明噪声完全淹没了信号，无法有效提取目标信号。从示波器的结果可以看出，当输入信噪比较高时（如 0 dB），示波器经过带通滤波处理可以对信号有所改善；但在较低信噪比（如 -30 dB 及以下）时，信号几乎被噪声完全掩盖，测量效果较差。

锁相放大器的测量通过控制量程灵敏度和时间常数来优化输出的信号。其测量结果显示出与示波器显著不同的表现。在 0 dB SNR 输入信号下，信号有效值为 120.1 mVrms，噪声有效值为 0.02 mVrms，输出信噪比为 75.57 dB。锁相放大器在输入信噪比为 0 dB 的情况下，显著提高了输出信噪比。在 -30 dB SNR 输入信号下，信号有效值为 3.77 mVrms，噪声有效值为 0.028 mVrms，输出信噪比为 42.58 dB。锁相放大器通过长时间常数和窄带滤波，有效抑制了噪声，使信噪比得到极大提升。在 -80 dB SNR 输入信号下，信号有效值为 12.1 mVrms，噪声有效值为 0.106 mVrms，输出信噪比为 41.15 dB。在如此低的信噪比下，锁相放大器依然能够将信号从噪声中提取出来，并保持较高的输出信噪比。锁相放大器在低信噪比条件下表现非常出色，主要得益于其窄带滤波器和同步检测功能，这使得即使信号被强噪声淹没，它仍然可以有效地提取信号。

示波器在高信噪比条件下能够较好地测量信号，并直观地显示波形，而在低信噪比情况下其测量效果较差。锁相放大器具有极高的信号提取能力，尤其在低信噪比情况下表现出色，其优势在于同步检测和窄带滤波，能够有效地从噪声中提取出特定频率的信号。然而，锁相放大器的响应速度相对较慢，尤其是选择长时间常数和窄带滤波时，信号的提取需要更长的时间，这使得它不适合需要实时观察信号变化的场合。此外，锁相放大器只能用于测量与参考信号频率相同的信号，因此应用范围比示波器要窄一些。

选择锁相放大器窄带宽的低通滤波器（LPF）可以有效降低噪声密度，从而提高测量信号的信噪比。在式 (D1-34) 中，噪声密度 \( n_{\omega_r} \) 表示在目标频率附近的噪声功率密度。低通滤波器的带宽 \( \Delta f \) 与噪声密度有关，噪声有效值 \( N \) 通常可以用以下公式表示：

\[
N = n_{\omega_r} \sqrt{\Delta f}
\]

从这个公式可以看出，带宽越窄，噪声有效值 \( N \) 就越小，因此信噪比会提高。锁相放大器通过选择较小的低通滤波器带宽，可以降低测量系统中的噪声水平，从而使信号在总输出中的占比增大，即提高输出信噪比。参考式 (D1-10)，信噪比 \( SNR \) 的提高与低通滤波器的带宽 \( \Delta f \) 有密切关系：

\[
SNR = \frac{R_{\text{signal}}}{N} \propto \frac{1}{\sqrt{\Delta f}}
\]

因此，降低低通滤波器的带宽可以有效提高信噪比。带宽减小使噪声密度乘积减小，从而降低噪声的有效值 \( N \)，提高了信噪比 \( SNR \)。这也是为什么锁相放大器在低信噪比情况下仍能获得较好测量效果的原因之一。

综上所述，示波器在高信噪比条件下能够较好地测量信号，并直观地显示波形，而在低信噪比情况下其测量效果较差；锁相放大器在低信噪比条件下依然可以通过同步检测和窄带滤波有效提取信号，其信噪比随着低通滤波器带宽的减小而显著提高，但相对测量速度较慢，适合特定频率信号的检测。因此，锁相放大器更适合需要在强噪声背景下检测微弱信号的应用，而示波器则更适合对信号波形、幅度进行实时监测的场合。
